\documentclass[12pt,letterpaper,boxed]{hmcpset}
\usepackage[margin=1in,headheight=14pt]{geometry}
\usepackage{amsfonts, amsmath, amssymb, enumerate, fancyhdr, gensymb, lastpage, mathtools, parskip, graphicx}
\usepackage{xcolor, tikz-cd}
\newcommand{\wg}[1]{\textcolor{violet}{#1}}
\newcommand{\OO}{\mathcal O}
\newcommand{\Q}{\mathbb Q}
\newcommand{\R}{\mathbb R}
\newcommand{\C}{\mathbb C}
\newcommand{\Z}{\mathbb Z}
%\newcommand{\H}{\mathbb H}
\newcommand{\abs}[1]{\left|#1\right|}
\newcommand{\im}{\text{im }}
\newcommand{\inv}{^{-1}}
\newcommand{\normal}{\unlhd} %% one can also use \trianglelelefteq
\newcommand{\anglee}[1]{\langle #1 \rangle}
\usepackage[shortlabels]{enumitem}

% Numbering macros
\pagestyle{fancy}
\lhead{Will Gilroy}
\chead{Algs Homework \#}
\rhead{03 November 2021}
\lfoot{}
\cfoot{}
\rfoot{Page\ \thepage\ of\ \pageref{LastPage}}

\linespread{1.5}

\newcommand\blankpage{
    \thispagestyle{empty}
    \addtocounter{page}{-1}
    \newpage}
\renewcommand\footrulewidth{0.4pt}

\begin{document}

\problemlist{Algorithms HW } 

%------------------------- Problem 1 -----------------------

\begin{problem}
	\includegraphics[scale=0.8]{1.png}
	\hfill
\end{problem}

\begin{solution}
We show that automorphisms of $\mathbb H$ are isometries by direct
calculation. Let $f \in PSL_2(\R)$, $z = x + iy \in \mathbb H$, and $w
= f(z)$. If we allow $\abs{dz}^2$ to denote $\abs{dz}^2 = dx^2 + dy^2$ then our
metric is $\frac{\abs{dz}^2}{y^2}$. To show that $f$ is an isometry we need
to show that 
\[
	\frac{\abs{dw}^2}{Im(f)^2} = \frac{\abs{dz}^2}{y^2}. 
\]
We have that $dw = f'(z) dz$ and so let us compute the derivative
$f'(z)$. Recall that $f(z) = \frac{az + b}{cz + d}$ for $a,b,c,d \in
\R$ and $ad-bc = 1$. Then using the quotient rule gives 
\begin{align*}
	f'(z) &= \frac{a}{cz+d} - \frac{c(dz+b)}{(cz+d)^2} \\
	&= \frac{acz + ad - acz - bc}{(cz+d)^2} \\
	&= \frac{1}{(cz+d)^2}.
\end{align*}
And so 
\[
	\frac{\abs{dw}^2}{\abs{dz}^2} = \frac{1}{\abs{cz+d}^2}.
\]
On the other hand let us compute $Im(f)^2$. 
We have 
\begin{align*}
	w &= \frac{az+b}{cz+d} \\
	&= \frac{(az+b)\overline{(cz+d)}}{\abs{cz+d}^2} \\
	&= \frac{(az+b)(c\bar z+d)}{\abs{cz+d}^2} && a,b,c,d \in \R \\
	&= \frac{ac\abs{z}^2 + adz + bc \bar z + bd}{\abs{cz+d}^2}.
\end{align*}
Now, notice that the denominator is a real number and that 
the only part of the numerator which may contain imaginary terms is
$adz + bc\bar z$. Consider 
\begin{align*}
	adz + bc\bar z &= ad(x+iy) + bc(x - iy) \\
	Im(adz + bc\bar z) &= ady - bcy \\
	&= y. 
\end{align*}
That is, 
\[
	\frac{Im(w)^2}{Im(z)^2} = \frac{1}{\abs{cz+d}^2} \frac{y^2}{y^2} = \frac{1}{\abs{cz+d}^2}. 
\]
We have thus shown that automorphisms of $\mathbb H$ are isometries of
the hyperbolic metric. 

Now let $f : \mathbb H \to \mathbb H$ be an isometry of the upper half
plane with the hyperbolic metric. 
We have that, $ds^2 = \frac{dx^2+dy^2}{y^2}$ which is a conformal
transform of the euclidean metric. Which means that its isometries
must be conformal. We have seen that the only conformal functions of
$\C$ are the holomorphic functions and the antiholomorphic functions.
An antiholomorphic function cannot be a function $\mathbb H \to
\mathbb H$. And so $f$ is a holomorphic function of $\mathbb H$. 

Now that we know $f$ is holomorphic, the calculation above shows that
$f$ must satisfy 
\[
	\frac{\abs{f'(z)}^2}{Im(f(z))} = \frac{1}{Im(z)}.
\]
If $f(x,y) = u(x,y) + iv(x,y)$ this expression gives us a condition on $v$. 

\end{solution}

\newpage

%------------------------- Problem 2 -----------------------

\begin{problem}
	\includegraphics[scale=0.8]{2.png}
	\hfill
\end{problem}

\begin{solution}
Let $p,q$ be two points of the upper half plane with the same real
part. If $\gamma$ is a curve in the upper half plane, the length of
$\gamma$ under the hyperbolic metric is given by contracting the
tangent vectors against the metric along the curve 
\[
??
\]
If $\gamma$ is not the vertical line connecting $p,q$ then it moves a
nonzero amount parallel to the real line. For a fixed imaginary value,
the hyperbolic length of a horizontal line is given by a scaling of
its euclidean length. In particular, any distance travelled in the
real direction adds a positive value to the length of our curve.
In particular this means that the geodesic connecting $p,q$ is given
by the curve which does not travel in the real direction at all ---
this is namely the vertical line connecting $p,q$.
\wg{Re-read this sentence to see if it makes any sense.}

Now suppose that $p,q$ two points no longer lying on a vertical line.
We saw in the previous question that any automorphism of $\mathbb H$
is a hyperbolic isometry. In particular, it must bring geodesics to
geodesics.
Then, it is sufficient to show that there
exists an automorphism of $\mathbb H$ which brings $p,q$ to a vertical
line. For a given $p,q$ there's a unique semicircle meeting the real
axis at a right angle specified by three points in the complex plane
$p,q,x$ where $x \in \R$. By triple transitivity of the Mobius transformations, there's
then a unique Mobius transformation $M$ mapping $(x,p,q) \mapsto (0,
\tilde p, \tilde q)$ where $\tilde p, \tilde q \in \mathbb H$ and are
colinear.

A priori this $M \in PGL_2(\C)$, however, if we can show that $M \in
PSL_2(\R)$ then are done. Since $M \in GL_2(\C)$ it has determinant $D
\neq 0$. We can scale $M$ by any $\lambda \in \C$ to acquire an
equivalent transformation. But $\abs{\lambda M} = \lambda^2 D$,
and so scaling $M$ by $1/\sqrt{D}$ means that $M$ can be considered to
have determinant $1$ without loss of generality.


Recall that the automorphisms of $\mathbb H$ are a subgroup of
$PSL_2(\R) \leq PGL_2(\C)$ and so, we need to argue that the
transformation found above has real entires and determinant $1$
\wg{Come back to this}


\end{solution}
\newpage



%------------------------- Problem 3 -----------------------

\begin{problem}[3]
	\hfill
\end{problem}
\begin{solution}
\end{solution}

\newpage

%------------------------- Problem 4 -----------------------

\begin{problem}
	\includegraphics[scale=0.8]{4.png}
	\hfill
\end{problem}
\begin{solution}
Let $f$ be as given in the question. First suppose that $f(0) = 0$
recall the Schwarz lemma for such a function states 
\[
	\abs{f(z)} \leq \abs{z} \text{for all } z \in \Delta 
	\qquad
	\text{ and } 
	\qquad \abs{f'(0)} \leq 1,
\]
and if equality is achieved in either for some $z_0 \in \Delta$, then $f$ is a rotation. 
Suppose $f$ has a fixed point, i.e. $f(z_0) = z_0$ for some $z_0 \in
\Delta$. But then we have 
$\abs{f(z_0)} = \abs{z_0}$ and so $f$ is a rotation. And so $f$ just
has the single fixed point at the origin. 

Now let $f$ be general (i.e. drop the assumption $f(0) = 0$). 
Since $f : \Delta \to \Delta$ we have that $f(0) = a$ for some $a \in
\Delta$. 
Now, recall that the Blaschke map $\phi_a$ is an automorphism of $\Delta$ which
exchanges $0$ and $a$. 
Now the map $h = \phi_a \circ f : \Delta \to \Delta$ is
holomorphic, as the composition of two holomorphic functions, and
fixes $0$. 
From the argument in the above paragraph can only have the origin as a
fixed point. 
\wg{What can we say from here?}

We have shown that if $f$ is non identity then it can only have a
single fixed point. The constant map $f: \Delta \to \Delta$ via $f(z) = a$ for
some $a \in \Delta$ has a fixed point at $a$.
If $f: \Delta \to \Delta$ is the identity restricted to the
disc has every point as a fixed point.

Overall the possibilities for the set of fixed points are $\Delta,
\{a\}$ where $a \in \Delta$ \wg{can we have no fixed points?}

\end{solution}


\newpage

%------------------------- Problem 6 -----------------------

\begin{problem}
	\includegraphics[scale=0.8]{6.png}
	\hfill
\end{problem}

\begin{solution}
First, we need to verify that the proposed $f$ has outputs to the
upper half plane. And then, if we can construct an inverse function,
then we have shown that the given function is a biholomorphism $\tilde
\Delta \to \mathbb H$ where $\tilde \Delta$ denotes the upper half disc. 

\wg{Come back and argue that the outputs of the function are in the upper half plane}

Let $w \in \mathbb H$. We want to show that there is a unique $z \in
\tilde \Delta$ such that $w = -1/2(z + 1/z)$. Rearranging this equation
gives 
\[
	0 = z^2 + 2wz + 1. 
\]
As a quadratic equation, this has two roots $z_0, z_1 \in \C$. 
Recalling then that we can factor our quadratic as $0 = (z -z_0)(z-z_1)$, 
the coefficients from this quadratic give the following relations on
the roots
\begin{align}
	z_0 z_1 &= 1 \\
	z_0 + z_1 &= -2w. 
\end{align}
This first equation gives us that $z_0$ and $z_1$ are inverses of each
other. I claim that neither $z_0, z_1$ lie on $S^1$. Suppose $z_0 \in
S^1$ then equation $1$ would imply that $z_1 \in S^1$ and that $z_0,
z_1$ are conjugate to each other. But then their sum is real. And then
equation $2$ cannot be satisfied since $w \in \mathbb H$ and so $-2w$
has strictly negative imaginary part. Hence, neither $z_0, z_1$ lie on
$S^1$. 

Now, equation $1$ implies that one of the roots lies in $\Delta$ and
the other root lies in $\mathbb C \setminus \Delta$, since $z \mapsto
1/z$ exchanges the unit circle and its complement. Let's say that $z_0
\in \Delta$. 
Now I claim that equation $2$ can only be satisifed if in fact $z_0
\in \tilde \Delta$. 
If we write $z_0 = r e^{i \theta}$ for $r < 1$ then $z_1 = z_0^{-1} =
1/r e^{-i \theta}$ (i.e. $z_1$ lies on the same ray as
$\overline{z_0}$ but with larger magnitude). And so if $z_0$ were to
lie in the lower half disc we would have that $z_0 + z_1$ has positive
imaginary part. But as noted earlier, $w \in \mathbb H$ and so $-2w$
has negative imaginary part. Thus, we must have $z_0 \in \tilde
\Delta$.


We have shown that given $w \in \mathbb H$ the equation $0 = z^2 + 2wz
+ 1$ has a unique solution in $\tilde \Delta$. In other words, 
we have shown that there exists a function $\mathbb H \to \tilde
\Delta$ which inverts the given $f$. Hence $f$ is a holomorphic
bijection, thus $f$ is a biholomorphism $\tilde \Delta \to \mathbb H$. 

\end{solution}

\newpage



\end{document}
